% Options for packages loaded elsewhere
\PassOptionsToPackage{unicode}{hyperref}
\PassOptionsToPackage{hyphens}{url}
%
\documentclass[
]{article}
\usepackage{amsmath,amssymb}
\usepackage{iftex}
\ifPDFTeX
  \usepackage[T1]{fontenc}
  \usepackage[utf8]{inputenc}
  \usepackage{textcomp} % provide euro and other symbols
\else % if luatex or xetex
  \usepackage{unicode-math} % this also loads fontspec
  \defaultfontfeatures{Scale=MatchLowercase}
  \defaultfontfeatures[\rmfamily]{Ligatures=TeX,Scale=1}
\fi
\usepackage{lmodern}
\ifPDFTeX\else
  % xetex/luatex font selection
\fi
% Use upquote if available, for straight quotes in verbatim environments
\IfFileExists{upquote.sty}{\usepackage{upquote}}{}
\IfFileExists{microtype.sty}{% use microtype if available
  \usepackage[]{microtype}
  \UseMicrotypeSet[protrusion]{basicmath} % disable protrusion for tt fonts
}{}
\makeatletter
\@ifundefined{KOMAClassName}{% if non-KOMA class
  \IfFileExists{parskip.sty}{%
    \usepackage{parskip}
  }{% else
    \setlength{\parindent}{0pt}
    \setlength{\parskip}{6pt plus 2pt minus 1pt}}
}{% if KOMA class
  \KOMAoptions{parskip=half}}
\makeatother
\usepackage{xcolor}
\usepackage{color}
\usepackage{fancyvrb}
\newcommand{\VerbBar}{|}
\newcommand{\VERB}{\Verb[commandchars=\\\{\}]}
\DefineVerbatimEnvironment{Highlighting}{Verbatim}{commandchars=\\\{\}}
% Add ',fontsize=\small' for more characters per line
\newenvironment{Shaded}{}{}
\newcommand{\AlertTok}[1]{\textcolor[rgb]{1.00,0.00,0.00}{\textbf{#1}}}
\newcommand{\AnnotationTok}[1]{\textcolor[rgb]{0.38,0.63,0.69}{\textbf{\textit{#1}}}}
\newcommand{\AttributeTok}[1]{\textcolor[rgb]{0.49,0.56,0.16}{#1}}
\newcommand{\BaseNTok}[1]{\textcolor[rgb]{0.25,0.63,0.44}{#1}}
\newcommand{\BuiltInTok}[1]{\textcolor[rgb]{0.00,0.50,0.00}{#1}}
\newcommand{\CharTok}[1]{\textcolor[rgb]{0.25,0.44,0.63}{#1}}
\newcommand{\CommentTok}[1]{\textcolor[rgb]{0.38,0.63,0.69}{\textit{#1}}}
\newcommand{\CommentVarTok}[1]{\textcolor[rgb]{0.38,0.63,0.69}{\textbf{\textit{#1}}}}
\newcommand{\ConstantTok}[1]{\textcolor[rgb]{0.53,0.00,0.00}{#1}}
\newcommand{\ControlFlowTok}[1]{\textcolor[rgb]{0.00,0.44,0.13}{\textbf{#1}}}
\newcommand{\DataTypeTok}[1]{\textcolor[rgb]{0.56,0.13,0.00}{#1}}
\newcommand{\DecValTok}[1]{\textcolor[rgb]{0.25,0.63,0.44}{#1}}
\newcommand{\DocumentationTok}[1]{\textcolor[rgb]{0.73,0.13,0.13}{\textit{#1}}}
\newcommand{\ErrorTok}[1]{\textcolor[rgb]{1.00,0.00,0.00}{\textbf{#1}}}
\newcommand{\ExtensionTok}[1]{#1}
\newcommand{\FloatTok}[1]{\textcolor[rgb]{0.25,0.63,0.44}{#1}}
\newcommand{\FunctionTok}[1]{\textcolor[rgb]{0.02,0.16,0.49}{#1}}
\newcommand{\ImportTok}[1]{\textcolor[rgb]{0.00,0.50,0.00}{\textbf{#1}}}
\newcommand{\InformationTok}[1]{\textcolor[rgb]{0.38,0.63,0.69}{\textbf{\textit{#1}}}}
\newcommand{\KeywordTok}[1]{\textcolor[rgb]{0.00,0.44,0.13}{\textbf{#1}}}
\newcommand{\NormalTok}[1]{#1}
\newcommand{\OperatorTok}[1]{\textcolor[rgb]{0.40,0.40,0.40}{#1}}
\newcommand{\OtherTok}[1]{\textcolor[rgb]{0.00,0.44,0.13}{#1}}
\newcommand{\PreprocessorTok}[1]{\textcolor[rgb]{0.74,0.48,0.00}{#1}}
\newcommand{\RegionMarkerTok}[1]{#1}
\newcommand{\SpecialCharTok}[1]{\textcolor[rgb]{0.25,0.44,0.63}{#1}}
\newcommand{\SpecialStringTok}[1]{\textcolor[rgb]{0.73,0.40,0.53}{#1}}
\newcommand{\StringTok}[1]{\textcolor[rgb]{0.25,0.44,0.63}{#1}}
\newcommand{\VariableTok}[1]{\textcolor[rgb]{0.10,0.09,0.49}{#1}}
\newcommand{\VerbatimStringTok}[1]{\textcolor[rgb]{0.25,0.44,0.63}{#1}}
\newcommand{\WarningTok}[1]{\textcolor[rgb]{0.38,0.63,0.69}{\textbf{\textit{#1}}}}
\usepackage{longtable,booktabs,array}
\usepackage{calc} % for calculating minipage widths
% Correct order of tables after \paragraph or \subparagraph
\usepackage{etoolbox}
\makeatletter
\patchcmd\longtable{\par}{\if@noskipsec\mbox{}\fi\par}{}{}
\makeatother
% Allow footnotes in longtable head/foot
\IfFileExists{footnotehyper.sty}{\usepackage{footnotehyper}}{\usepackage{footnote}}
\makesavenoteenv{longtable}
\setlength{\emergencystretch}{3em} % prevent overfull lines
\providecommand{\tightlist}{%
  \setlength{\itemsep}{0pt}\setlength{\parskip}{0pt}}
\setcounter{secnumdepth}{-\maxdimen} % remove section numbering
\newlength{\cslhangindent}
\setlength{\cslhangindent}{1.5em}
\newlength{\csllabelwidth}
\setlength{\csllabelwidth}{3em}
\newlength{\cslentryspacingunit} % times entry-spacing
\setlength{\cslentryspacingunit}{\parskip}
\newenvironment{CSLReferences}[2] % #1 hanging-ident, #2 entry spacing
 {% don't indent paragraphs
  \setlength{\parindent}{0pt}
  % turn on hanging indent if param 1 is 1
  \ifodd #1
  \let\oldpar\par
  \def\par{\hangindent=\cslhangindent\oldpar}
  \fi
  % set entry spacing
  \setlength{\parskip}{#2\cslentryspacingunit}
 }%
 {}
\usepackage{calc}
\newcommand{\CSLBlock}[1]{#1\hfill\break}
\newcommand{\CSLLeftMargin}[1]{\parbox[t]{\csllabelwidth}{#1}}
\newcommand{\CSLRightInline}[1]{\parbox[t]{\linewidth - \csllabelwidth}{#1}\break}
\newcommand{\CSLIndent}[1]{\hspace{\cslhangindent}#1}
\ifLuaTeX
  \usepackage{selnolig}  % disable illegal ligatures
\fi
\IfFileExists{bookmark.sty}{\usepackage{bookmark}}{\usepackage{hyperref}}
\IfFileExists{xurl.sty}{\usepackage{xurl}}{} % add URL line breaks if available
\urlstyle{same}
\hypersetup{
  pdftitle={When Low WER Becomes Dangerous: Counterfactual Decision-Stability ASR for High-Stakes Speech-Driven Decision Systems},
  hidelinks,
  pdfcreator={LaTeX via pandoc}}

\title{When Low WER Becomes Dangerous: Counterfactual Decision-Stability
ASR for High-Stakes Speech-Driven Decision Systems}
\author{}
\date{2026-05-28}

\begin{document}
\maketitle

\hypertarget{manuscript-draft}{%
\section{Manuscript Draft}\label{manuscript-draft}}

Working title: When Low WER Becomes Dangerous: Counterfactual
Decision-Stability ASR for High-Stakes Speech-Driven Decision Systems

Status: paper-facing v0 draft with second-reviewer evidence-boundary
revisions, results table package, reproducibility layer, scope control,
limitations, and artifact availability plan.

Date: 2026-05-28

\hypertarget{evidence-boundary-freeze}{%
\subsection{Evidence Boundary Freeze}\label{evidence-boundary-freeze}}

This manuscript uses a scoped evidence chain:

\begin{enumerate}
\def\labelenumi{\arabic{enumi}.}
\tightlist
\item
  The 258-row test split is used as scope-controlled split and
  model-comparison evidence.
\item
  The selected-300 proxy outputs are used as input provenance and
  row-selection evidence.
\item
  The selected-300 human-reviewed predictor and human-reviewed recovery
  outputs are the paper-grade risk and recovery evidence.
\end{enumerate}

Reference transcripts remain accepted human-reviewed ground truth for
WER/CER scoring. Transcript review is not reopened unless a future
review task changes the field scope or challenges the accepted reference
transcript content.

Tracked manuscript artifacts must remain aggregate-only. Raw audio, raw
transcripts, audio IDs, selected sample IDs, model hypotheses, local
response sheets, reviewer notes, runtime logs, raw predictions, and
model weights remain local or ignored and are not committed.

Authoritative boundary records:

\begin{itemize}
\tightlist
\item
  \texttt{70\_experiments/runs/postdoc\_evidence\_chain\_2026\_05\_25/evidence\_chain\_readiness\_summary.json}
\item
  \texttt{70\_experiments/runs/postdoc\_evidence\_chain\_2026\_05\_25/publishable\_evidence\_completion\_summary.json}
\item
  \texttt{70\_experiments/runs/postdoc\_evidence\_chain\_2026\_05\_25/consequence\_evidence\_matrix.tsv}
\item
  \texttt{70\_experiments/runs/janus\_300\_high\_stakes\_human\_audit\_selection\_2026\_05\_25/human\_audit\_refresh\_summary.json}
\end{itemize}

\hypertarget{abstract}{%
\subsection{Abstract}\label{abstract}}

Speech-driven decision systems increasingly transform conversations into
operational signals for routing, escalation, compliance monitoring, and
case-handling decisions. In high-stakes anti-fraud calls, small ASR
differences can land on decision atoms such as negation, amount, actor,
action, time, and intent. A transcript can therefore remain close under
WER or CER while a plausible alternative transcript changes the
downstream decision.

We propose Counterfactual Decision-Stability ASR (CDS-ASR), a framework
that evaluates whether downstream decisions remain stable under
acoustically and semantically plausible transcript alternatives. CDS-ASR
extracts risk atoms, constructs plausible ASR alternatives, scores
decision instability with Counterfactual Escalation Instability Score
(CEIS), and applies constrained recovery or conservative machine action
for high-risk spans while preserving aggregate-only reproducibility for
sensitive call evidence.

Across a six-model 258-row split comparison, selected-300 high-stakes
provenance outputs, and a 30-row/90-model-assessment human-reviewed
audit, CDS-ASR provides a consequence-centered evidence layer beyond
transcript similarity. In the scoped selected-300 human-reviewed audit,
CEIS achieves the highest decision-change AUC and a zero-false-negative
operating point. In aggregate policy replay, risk-triggered policies,
including CEIS-triggered conservative action, eliminate high-risk missed
and critical miss counts under the aggregate-only claim boundary. These
results support decision-stability evaluation as a scoped companion to
transcript accuracy, semantic metrics, and confidence-aware correction.

\hypertarget{introduction}{%
\subsection{Introduction}\label{introduction}}

Contact-center speech is becoming an operational input rather than a
passive record. Commercial analytics systems already turn voice
conversations and transcripts into categories, summaries, alerts,
compliance cues, and routing signals (Amazon Web Services 2026a, 2026b).
This workflow makes ASR output part of operational evidence: a
transcript can support categorization, summarization, compliance review,
alerts, and routing before a human ever reads the full call record.

Anti-fraud hotlines provide a concrete high-stakes setting. Taiwan's
National Police Agency describes the 165 hotline as a public anti-fraud
channel where staff record incident details and provide information to
victims (National Police Agency, Ministry of the Interior 2024). The
scale of cyber-enabled fraud makes this triage problem consequential:
the FBI's 2025 Internet Crime Report context reported more than one
million IC3 complaints and large cyber-enabled fraud losses (Federal
Bureau of Investigation 2026b, 2026a). In this setting, safer
speech-driven triage is not only a transcription problem; it is a
decision problem. The risk is concentrated in small spoken details:
whether the caller already transferred money, whether the amount was
30,000 or 300,000, who initiated contact, when the event happened, and
whether the caller is certain. These are decision atoms. When ASR
changes a decision atom, the downstream action can change even if most
words remain similar.

In this paper, ``dangerous'' is operationalized as ASR-induced decision
change that can produce unsafe downrouting, high-risk missed escalation,
or critical miss under the declared anti-fraud triage policy. This
definition makes the title a measurable evidence claim rather than a
general safety label.

Current ASR evaluation and repair methods provide strong foundations for
this setting. WER and CER give reproducible transcript-centered
baselines. Semantic ASR metrics make evaluation more task-aware: Kim et
al.~show that WER can fail to reflect downstream
natural-language-understanding behavior, and Rugayan et al.~evaluate
Aligned Semantic Distance as a semantic metric with stronger alignment
to human judgments and downstream tasks (Kim et al. 2021; Rugayan,
Salvi, and Svendsen 2023).

LLM-based and confidence-aware correction methods strengthen the
transcript repair layer. Naderi et al.~study post-hoc ASR correction
with LLMs and use confidence-based filtering to reduce harmful changes
to likely accurate transcripts (Naderi et al. 2024). Selective
prediction and reject-option work also show how machine systems can
trade coverage for lower error in settings where acting on uncertain
predictions is costly (Chow 1970; Geifman and El-Yaniv 2017, 2019;
Angelopoulos and Bates 2021). These methods make the evidence chain more
informative, and they motivate the next question for high-stakes speech
workflows: whether plausible ASR alternatives preserve the downstream
decision.

This decision-stability gap appears when a low edit-distance transcript
still leaves uncertainty around negation, amount, action, actor, time,
intent, or scam-pattern atoms. A transcript can remain close under
WER/CER and still reverse whether a case should remain in routine
review, move to priority review, or trigger critical escalation. The
central research question is therefore decision-centered: would a
plausible ASR alternative change the downstream high-stakes decision?

We propose Counterfactual Decision-Stability ASR (CDS-ASR) to answer
that question directly. The framework extracts risk atoms, generates
acoustically and semantically plausible transcript alternatives,
measures decision instability with CEIS, and applies constrained
recovery or conservative machine action for high-risk spans. WER, CER,
semantic metrics, SRES, confidence, and correction baselines remain
comparison layers; the central contribution is the decision-stability
test that connects transcript uncertainty back to the opening real-world
decision problem.

Verified source anchors for this introduction are maintained in
\texttt{citation\_seed.md} and \texttt{references.bib}, with access
dates refreshed on 2026-05-28. Empirical manuscript claims point to the
aggregate artifacts listed in the Appendix.

\hypertarget{related-work}{%
\subsection{Related Work}\label{related-work}}

\hypertarget{transcript-centered-asr-metrics}{%
\subsubsection{Transcript-Centered ASR
Metrics}\label{transcript-centered-asr-metrics}}

Transcript accuracy metrics provide the reproducible baseline for ASR
comparison. WER and CER remain necessary because they allow stable
comparison across models and splits when tokenization, normalization,
micro/macro scope, and zero-reference handling are declared. This
manuscript keeps WER/CER as auditable surface metrics and uses the
repo's metric audit to define the Chinese ASR reporting policy. The
paper does not discard transcript accuracy; it treats WER/CER as the
auditable surface layer that must be reported before stronger downstream
claims are made.

\hypertarget{semantic-and-downstream-aware-asr-evaluation}{%
\subsubsection{Semantic And Downstream-Aware ASR
Evaluation}\label{semantic-and-downstream-aware-asr-evaluation}}

Semantic ASR metrics extend transcript evaluation toward meaning
preservation and downstream task behavior. Kim et al.~introduce Semantic
Distance for ASR performance analysis toward spoken-language
understanding, motivating the move from literal correctness to semantic
correctness for intent recognition, slot filling, semantic parsing, and
named entity recognition (Kim et al. 2021). Rugayan et al.~evaluate
Aligned Semantic Distance against WER and show its value for human
perception and downstream NLP tasks (Rugayan, Salvi, and Svendsen 2023).
CDS-ASR uses this line of work as a foundation and adds a direct
decision-stability target: whether plausible ASR alternatives change a
high-stakes downstream decision.

\hypertarget{transcript-repair-and-conservative-decision-action}{%
\subsubsection{Transcript Repair And Conservative Decision
Action}\label{transcript-repair-and-conservative-decision-action}}

LLM and confidence-aware correction methods provide transcript-repair
baselines. Naderi et al.~show that confidence-based filtering can guide
post-hoc LLM correction of ASR transcripts and reduce the chance of
changing likely accurate transcripts (Naderi et al. 2024). CDS-ASR
complements this line by evaluating the decision interval that remains
when plausible ASR alternatives affect risk atoms. Recovery in this
manuscript is automatic and machine-bounded: human review supplies
evaluation labels and governance evidence, while risk-triggered policies
supply the scoped intervention test.

The conservative-action layer also connects to selective prediction and
reject-option research. Chow formalizes the error-reject tradeoff for
recognition systems (Chow 1970), while modern selective classification
and SelectiveNet show how neural systems can abstain or reject to manage
risk-coverage tradeoffs (Geifman and El-Yaniv 2017, 2019). Conformal
prediction further motivates distribution-free uncertainty sets for
high-risk systems (Angelopoulos and Bates 2021). CDS-ASR applies this
governance intuition to speech-driven decision evidence: when transcript
alternatives affect decision atoms, the system should expose the
instability or take conservative action rather than treat one transcript
as fully stable.

High-stakes speech domains reinforce the need for consequence-aware ASR
evaluation. In psychotherapy, Miner et al.~show that ASR feasibility
depends on the clinical use case and that individual-level safety
monitoring requires a more cautious threshold than population-level
language analysis (Miner et al. 2020). CDS-ASR brings the same
claim-evidence discipline to anti-fraud calls by separating transcript
accuracy, semantic risk, decision instability, and aggregate recovery
evidence.

This paper positions ASR models as evidence-producing components. The
protagonist is the downstream decision that changes when an acoustically
plausible transcript difference lands on a decision-critical atom.

\hypertarget{method}{%
\subsection{Method}\label{method}}

\hypertarget{problem-formulation}{%
\subsubsection{Problem Formulation}\label{problem-formulation}}

We study speech-driven decision systems where an audio segment is
transcribed by ASR and then used by a downstream escalation workflow.
The downstream label space includes routine review and higher-priority
escalation states. The central unit is a model-sample: one ASR
hypothesis for one audio row, evaluated against aggregate risk and
decision-stability evidence.

The paper asks whether the decision remains stable under plausible
transcript alternatives. This makes transcript accuracy a supporting
layer and decision-stability the primary high-stakes safety target.

\hypertarget{cds-asr-pipeline}{%
\subsubsection{CDS-ASR Pipeline}\label{cds-asr-pipeline}}

CDS-ASR evaluates high-stakes ASR through a decision-stability pipeline:

\begin{Shaded}
\begin{Highlighting}[]
\NormalTok{audio}
\NormalTok{{-}\textgreater{} ASR transcript and runtime/quality signals}
\NormalTok{{-}\textgreater{} risk atom extraction}
\NormalTok{{-}\textgreater{} plausible counterfactual transcript variants}
\NormalTok{{-}\textgreater{} downstream decision model}
\NormalTok{{-}\textgreater{} SRES and CEIS scoring}
\NormalTok{{-}\textgreater{} constrained recovery or conservative machine action}
\end{Highlighting}
\end{Shaded}

Risk atoms are transcript spans whose alternatives can affect downstream
decisions. This manuscript focuses on negation, amount, action, actor,
time, intent, uncertainty, and scam-pattern atoms.

CEIS scores the maximum decision-flip risk among plausible variants. In
this paper-facing implementation, the plausibility term is a bounded
proxy rather than a calibrated acoustic posterior:

\begin{Shaded}
\begin{Highlighting}[]
\NormalTok{CEIS(x) = max over v in V(x) [}
\NormalTok{    Plausibility(v | x) * RiskAtomWeight(v) * DecisionDistance(f(x), f(v))}
\NormalTok{]}
\end{Highlighting}
\end{Shaded}

\texttt{Plausibility(v\ \textbar{}\ x)} denotes a bounded proxy
plausibility score derived from model disagreement, Mandarin phonetic
ambiguity, domain-slot alternatives, and available ASR/runtime signals.
\texttt{RiskAtomWeight(v)} is a risk-atom class weight defined by the
decision-critical atom schema. \texttt{DecisionDistance} maps the
downstream label space into an ordinal safety distance over
\texttt{no\_escalation}, \texttt{review}, \texttt{priority\_review},
\texttt{critical\_escalation}, conservative machine action, and
abstention, with critical misses treated as the highest-risk direction
for conservative action.

\hypertarget{downstream-decision-function}{%
\subsubsection{Downstream Decision
Function}\label{downstream-decision-function}}

The downstream decision function \texttt{f} is fixed before CEIS
scoring. It maps each transcript or plausible variant into a declared
triage action:

\begin{Shaded}
\begin{Highlighting}[]
\NormalTok{f(transcript) in \{}
\NormalTok{    no\_escalation,}
\NormalTok{    manual\_review,}
\NormalTok{    priority\_review,}
\NormalTok{    critical\_escalation,}
\NormalTok{    conservative\_machine\_action,}
\NormalTok{    abstain}
\NormalTok{\}}
\end{Highlighting}
\end{Shaded}

The distance function is policy-aligned rather than
language-model-generated: same-action pairs have distance 0, neighboring
review/escalation states have distance 1, and movement from no
escalation or routine handling to critical escalation receives the
largest distance. Unsafe downrouting of a critical event receives the
maximum penalty used by the policy replay. The method contract is
versioned in \texttt{docs/downstream\_decision\_contract.md} and
\texttt{80\_semantic\_risk\_asr/scoring/ceis\_config.json} so CEIS
rankings can be reconstructed from aggregate artifacts.

\hypertarget{ceis-scale-discipline}{%
\subsubsection{CEIS Scale Discipline}\label{ceis-scale-discipline}}

All CEIS components are treated as bounded comparison terms, not
calibrated probabilities. \texttt{Plausibility(v\ \textbar{}\ x)} is a
proxy plausibility score; the paper does not claim an acoustic posterior
unless a later implementation supplies one. \texttt{RiskAtomWeight(v)}
and \texttt{DecisionDistance(f(x),\ f(v))} must be normalized or
tabulated before final submission. The next method-hardening analyses
are: uniform-weight ablation, no-plausibility ablation,
binary-decision-flip ablation, max versus top-k mean aggregation, and
CEIS behavior by atom class. Those analyses can be added only from
aggregate-safe recomputation or reported as planned validation if the
current artifact package cannot reconstruct them.

Thresholds reported in the current predictor table are diagnostic
operating points selected on the scoped reviewed audit. They are useful
for aggregate comparison and reviewer inspection, but they are not
frozen deployment thresholds unless later selected on a separate
development set.

Recovery remains automatic and machine-bounded. Human review is used as
an evaluation and governance layer, not as the proposed recovery method.
The recovery policies compare no recovery, confidence-only trigger,
SRES-triggered recovery, CEIS-triggered conservative action, and CEIS
ensemble arbitration.

\hypertarget{risk-atom-schema}{%
\subsubsection{Risk Atom Schema}\label{risk-atom-schema}}

The risk atom schema focuses on decision-critical spans:

\begin{longtable}[]{@{}
  >{\raggedright\arraybackslash}p{(\columnwidth - 4\tabcolsep) * \real{0.3333}}
  >{\raggedright\arraybackslash}p{(\columnwidth - 4\tabcolsep) * \real{0.3333}}
  >{\raggedright\arraybackslash}p{(\columnwidth - 4\tabcolsep) * \real{0.3333}}@{}}
\toprule\noalign{}
\begin{minipage}[b]{\linewidth}\raggedright
Risk atom
\end{minipage} & \begin{minipage}[b]{\linewidth}\raggedright
Decision role
\end{minipage} & \begin{minipage}[b]{\linewidth}\raggedright
Example effect
\end{minipage} \\
\midrule\noalign{}
\endhead
\bottomrule\noalign{}
\endlastfoot
Negation & reverses event state & transfer vs no transfer \\
Amount & changes severity & 30,000 vs 300,000 \\
Action & changes case stage & asking vs reporting vs transferring \\
Actor & changes scam interpretation & bank, police, family, customer
service \\
Time & changes urgency & today, yesterday, next week \\
Intent & changes routing & inquiry, report, complaint \\
Uncertainty & changes confidence and safe action & sure, maybe, not
sure \\
Scam pattern & changes case type & investment, fake police, installment
cancellation \\
\end{longtable}

\hypertarget{counterfactual-variant-contract}{%
\subsubsection{Counterfactual Variant
Contract}\label{counterfactual-variant-contract}}

Counterfactual variants are plausible ASR alternatives concentrated
around risk atoms. They are not generic paraphrases. A variant can be
supported by acoustic ambiguity, Mandarin phonetic confusion,
domain-slot alternatives, or model disagreement. The paper-facing claim
is not that every possible variant is enumerated; it is that the
generated alternatives make downstream decision-instability measurable.

\hypertarget{recovery-policy-contract}{%
\subsubsection{Recovery Policy
Contract}\label{recovery-policy-contract}}

The five-policy recovery experiment is:

\begin{enumerate}
\def\labelenumi{\arabic{enumi}.}
\tightlist
\item
  no recovery;
\item
  confidence-only trigger;
\item
  SRES-triggered recovery;
\item
  CEIS-triggered conservative action;
\item
  CEIS ensemble arbitration.
\end{enumerate}

The scoped recovery claim is evaluated only on aggregate human-reviewed
selected-300 evidence. Per-row details stay local or ignored. At the
policy layer, both SRES-triggered recovery and CEIS-triggered
conservative action are reported as risk-triggered conservative
policies; CEIS's distinct contribution is evaluated primarily in the
predictor layer and in its decision-stability framing.

\hypertarget{reproducibility-layer}{%
\subsubsection{Reproducibility Layer}\label{reproducibility-layer}}

The paper-facing ASR surface metric is \texttt{cer\_zh\_micro}. The
supplemental word metric is \texttt{wer\_zh\_jieba\_micro}. Raw
whitespace WER and legacy stored WER fields are audit-only because they
are unstable for unsegmented Chinese transcripts.

The metric policy is validated by
\texttt{70\_experiments/runs/wer\_metric\_audit\_2026\_05\_25/}, which
records manifest checks, zero-reference-unit checks, tokenizer and
normalization policy, and \texttt{jiwer} cross-checks.

The target transcription locale is Taiwan Traditional Chinese. New
candidate models must satisfy the strict Taiwan Traditional Chinese
locale gate before promotion. Previously completed comparable baselines
are retained only as disclosed baselines, with locale-violation counts
reported and no promotion claim attached. Post-decode OpenCC or other
conversion can be evaluated only as a deployment repair lane; it cannot
be used to claim that the raw ASR model passed the locale gate.

Selected-300 human audit evidence is tracked only through aggregate
outputs. The local transcript-bearing audit sheet remains ignored. The
reviewed scope is risk atoms, decision-change labels, expected safe
action, confidence, per-model assessment fields, and per-row timing. The
current aggregate status is \texttt{review\_complete} with 30/30
reviewed rows and 90/90 reviewed model assessments.

All paper-facing claims point to aggregate artifacts. The manuscript
does not require raw audio, raw transcripts, selected row IDs,
hypotheses, reviewer notes, or model weights.

\hypertarget{evaluation-units-and-n-ladder}{%
\subsubsection{Evaluation Units And
N-Ladder}\label{evaluation-units-and-n-ladder}}

The manuscript separates evidence units explicitly so that model
assessments are not mistaken for independent audio rows.

\begin{longtable}[]{@{}
  >{\raggedright\arraybackslash}p{(\columnwidth - 6\tabcolsep) * \real{0.2143}}
  >{\raggedleft\arraybackslash}p{(\columnwidth - 6\tabcolsep) * \real{0.2857}}
  >{\raggedleft\arraybackslash}p{(\columnwidth - 6\tabcolsep) * \real{0.2857}}
  >{\raggedright\arraybackslash}p{(\columnwidth - 6\tabcolsep) * \real{0.2143}}@{}}
\toprule\noalign{}
\begin{minipage}[b]{\linewidth}\raggedright
Layer
\end{minipage} & \begin{minipage}[b]{\linewidth}\raggedleft
Unit
\end{minipage} & \begin{minipage}[b]{\linewidth}\raggedleft
N
\end{minipage} & \begin{minipage}[b]{\linewidth}\raggedright
Role
\end{minipage} \\
\midrule\noalign{}
\endhead
\bottomrule\noalign{}
\endlastfoot
Test split & audio rows & 258 & ASR model comparison \\
Selected high-stakes provenance & selected candidate rows / outputs &
300 & row-selection and provenance \\
Human-reviewed audit rows & audio rows & 30 & decision-critical review
unit \\
Human-reviewed model assessments & model-row assessments & 90 &
predictor and recovery evaluation \\
\end{longtable}

Because the 90 model assessments are clustered within 30 audio rows,
inferential uncertainty should be estimated with row-clustered bootstrap
resampling or leave-one-row-out sensitivity analysis. Point estimates
are reported as scoped descriptive evidence; deployment claims do not
treat the 90 assessments as independent calls.

\hypertarget{selection-provenance-and-enrichment}{%
\subsubsection{Selection Provenance And
Enrichment}\label{selection-provenance-and-enrichment}}

The selected-300 audit surface is an enriched high-stakes sample, not a
prevalence-preserving sample of all anti-fraud calls. The selection
summary uses aggregate-safe provenance signals from SRES scoring, CEIS
scoring, and downstream escalation decisions, including high-risk
missed, critical miss, unsafe downrouting, model disagreement, high
proxy risk, risk-score fill, and clean-control strata. The aggregate
selection thresholds recorded for the current package are
\texttt{low\_wer\_threshold=10.0}, \texttt{sres\_threshold=20.0}, and
\texttt{ceis\_threshold=5.0}.

Predictor results therefore support metric behavior on an enriched
high-stakes audit surface rather than population-level risk prevalence.
If a submission emphasizes CEIS/SRES separation, a sensitivity check
should exclude rows directly selected by CEIS or SRES family signals and
confirm whether the same directional pattern remains.

\hypertarget{variant-coverage-and-human-review-reliability}{%
\subsubsection{Variant Coverage And Human Review
Reliability}\label{variant-coverage-and-human-review-reliability}}

The current aggregate-only coverage audit reports CEIS top-atom proxy
coverage over the 90 reviewed model assessments. It records 90 aggregate
proxy observations across 30 reviewed rows, with top-atom counts of
negation 47, amount 37, action 3, and actor 3. The reviewed selection
surface also covers risk-signal atoms in aggregate: action 25, actor 15,
amount 23, negation 14, and scam pattern 23 selected rows. These counts
support submission-safe coverage auditing, not release of variant text.
Source-specific coverage is available for model-disagreement provenance
in aggregate; phonetic, domain-slot, runtime-signal, rejected-variant,
and full variant-generation logs are not currently reconstructed as
release artifacts. Stronger generator claims therefore require a future
aggregate variant-generation log.

Human labels come from a completed expert audit rather than
multi-annotator adjudication. A second-reviewer blinded spot-check is
not included in the current submission-prep package, so inter-annotator
agreement is not claimed. The claim boundary is therefore a completed
single-expert audit over 30 reviewed rows and 90 model assessments. A
future validation extension can add a blinded 5-10 row spot-check while
preserving the frozen selected-300 review boundary.

\hypertarget{experiments}{%
\subsection{Experiments}\label{experiments}}

Experiment 1 evaluates comparable ASR model evidence on the canonical
258-row test split. This table supports model-comparison and split
evidence, not the paper-grade selected-300 risk/recovery claims.

Experiment 2 uses selected-300 high-stakes proxy outputs as input
provenance. The selected-300 proxy outputs identify an enriched
high-stakes evidence surface and the review sample, but their raw rows
and transcript-bearing metric inputs remain local or ignored. This
experiment supports selection provenance, not population prevalence.

Experiment 3 evaluates WER, CER, SRES, and CEIS against human-reviewed
decision-change labels over 90 model assessments from the selected-300
audit. These assessments are clustered within 30 reviewed audio rows.
This is the paper-grade predictor evidence for metric-insufficiency and
decision-stability claims, reported as scoped descriptive evidence until
row-clustered uncertainty is added.

Experiment 4 evaluates five recovery policies against the same
human-reviewed evidence layer. This is aggregate policy replay evidence,
not a live causal deployment trial.

Candidate ASR and multimodal models are kept in a separate exploratory
lane. They must not enter the main ASR table until they pass
field-contract, runtime-validity, and Taiwan Traditional Chinese locale
gates in order.

\hypertarget{results}{%
\subsection{Results}\label{results}}

The Results section separates comparable ASR evidence, candidate-lane
boundaries, human-reviewed predictor evidence, and human-reviewed
recovery evidence. This separation preserves claim-evidence alignment:
the main ASR benchmark supports model-comparison context, while
selected-300 human-reviewed outputs support paper-grade risk and
recovery claims.

Human review supplies evaluation labels only. The proposed recovery
policies are automatic, aggregate-evaluated policies; no
transcript-bearing row content is required for paper-facing claims.

\hypertarget{table-1.-main-asr-benchmark-table}{%
\subsubsection{Table 1. Main ASR Benchmark
Table}\label{table-1.-main-asr-benchmark-table}}

Table 1 reports six comparable ASR runs on the canonical 258-row split.
The paper-facing primary ASR metric is \texttt{cer\_zh\_micro};
\texttt{wer\_zh\_jieba\_micro} is reported as the supplemental segmented
word metric. Unsafe downrouting and high-risk missed counts are
scope-controlled split/model-comparison evidence, not the final
selected-300 human-reviewed risk claim.

Only models with completed comparable 258-row split evidence are
included.

\begin{longtable}[]{@{}
  >{\raggedright\arraybackslash}p{(\columnwidth - 16\tabcolsep) * \real{0.0909}}
  >{\raggedright\arraybackslash}p{(\columnwidth - 16\tabcolsep) * \real{0.0909}}
  >{\raggedleft\arraybackslash}p{(\columnwidth - 16\tabcolsep) * \real{0.1212}}
  >{\raggedleft\arraybackslash}p{(\columnwidth - 16\tabcolsep) * \real{0.1212}}
  >{\raggedleft\arraybackslash}p{(\columnwidth - 16\tabcolsep) * \real{0.1212}}
  >{\raggedleft\arraybackslash}p{(\columnwidth - 16\tabcolsep) * \real{0.1212}}
  >{\raggedleft\arraybackslash}p{(\columnwidth - 16\tabcolsep) * \real{0.1212}}
  >{\raggedleft\arraybackslash}p{(\columnwidth - 16\tabcolsep) * \real{0.1212}}
  >{\raggedright\arraybackslash}p{(\columnwidth - 16\tabcolsep) * \real{0.0909}}@{}}
\toprule\noalign{}
\begin{minipage}[b]{\linewidth}\raggedright
Run
\end{minipage} & \begin{minipage}[b]{\linewidth}\raggedright
Model family
\end{minipage} & \begin{minipage}[b]{\linewidth}\raggedleft
Rows
\end{minipage} & \begin{minipage}[b]{\linewidth}\raggedleft
\texttt{cer\_zh\_micro}
\end{minipage} & \begin{minipage}[b]{\linewidth}\raggedleft
\texttt{wer\_zh\_jieba\_micro}
\end{minipage} & \begin{minipage}[b]{\linewidth}\raggedleft
Unsafe downrouting
\end{minipage} & \begin{minipage}[b]{\linewidth}\raggedleft
High-risk missed
\end{minipage} & \begin{minipage}[b]{\linewidth}\raggedleft
Locale violation rows
\end{minipage} & \begin{minipage}[b]{\linewidth}\raggedright
Paper use
\end{minipage} \\
\midrule\noalign{}
\endhead
\bottomrule\noalign{}
\endlastfoot
\texttt{breeze\_asr25\_partial\_encoder\_legacy\_best\_test\_split} &
Breeze-ASR-25 partial encoder & 258 & 15.04 & 21.53 & 7 & 4 & 0 &
Scope-controlled split/model comparison \\
\texttt{breeze\_asr25\_lora\_legacy\_best\_test\_split} & Breeze-ASR-25
LoRA & 258 & 18.23 & 25.59 & 10 & 7 & 0 & Scope-controlled split/model
comparison \\
\texttt{breeze\_asr25\_base\_test\_split} & Breeze-ASR-25 base & 258 &
22.72 & 30.39 & 34 & 30 & 0 & Scope-controlled split/model comparison \\
\texttt{breeze\_asr26\_test\_split} & Breeze-ASR-26 & 258 & 24.27 &
32.29 & 27 & 22 & 0 & Optional dialect-aware comparator \\
\texttt{whisper\_large\_v2\_test\_split} & Whisper large-v2 & 258 &
24.72 & 32.23 & 33 & 28 & 1 & Comparable ASR baseline \\
\texttt{whisper\_small\_test\_split} & Whisper small & 258 & 34.86 &
43.44 & 76 & 70 & 4 & Comparable ASR baseline \\
\end{longtable}

Evidence source:
\texttt{70\_experiments/runs/janus\_258\_test\_split\_asr\_cds\_proxy/asr\_cds\_proxy\_comparison.tsv}.

The legacy Breeze-ASR-25 partial encoder is the strongest current
hypothesis generator on the 258-row split, with the lowest
\texttt{cer\_zh\_micro} and the lowest unsafe downrouting and high-risk
missed counts among the six comparable runs. This supports using it as a
strong ASR evidence layer while keeping final risk and recovery claims
tied to the human-reviewed selected-300 outputs.

\hypertarget{table-2.-candidate-and-exploratory-lane}{%
\subsubsection{Table 2. Candidate And Exploratory
Lane}\label{table-2.-candidate-and-exploratory-lane}}

Table 2 reports bounded candidate-lane evidence for models that are not
promoted to the main 258-row or selected-300 tables. These gates
document feasibility, locale behavior, and runtime constraints without
mixing exploratory candidates with comparable ASR baselines.

These models are bounded negative, feasibility, or runtime evidence.
They do not enter the main ASR benchmark table.

\begin{longtable}[]{@{}
  >{\raggedright\arraybackslash}p{(\columnwidth - 12\tabcolsep) * \real{0.1250}}
  >{\raggedright\arraybackslash}p{(\columnwidth - 12\tabcolsep) * \real{0.1250}}
  >{\raggedleft\arraybackslash}p{(\columnwidth - 12\tabcolsep) * \real{0.1667}}
  >{\raggedleft\arraybackslash}p{(\columnwidth - 12\tabcolsep) * \real{0.1667}}
  >{\raggedleft\arraybackslash}p{(\columnwidth - 12\tabcolsep) * \real{0.1667}}
  >{\raggedright\arraybackslash}p{(\columnwidth - 12\tabcolsep) * \real{0.1250}}
  >{\raggedright\arraybackslash}p{(\columnwidth - 12\tabcolsep) * \real{0.1250}}@{}}
\toprule\noalign{}
\begin{minipage}[b]{\linewidth}\raggedright
Candidate
\end{minipage} & \begin{minipage}[b]{\linewidth}\raggedright
Current gate
\end{minipage} & \begin{minipage}[b]{\linewidth}\raggedleft
Rows
\end{minipage} & \begin{minipage}[b]{\linewidth}\raggedleft
\texttt{cer\_zh\_micro}
\end{minipage} & \begin{minipage}[b]{\linewidth}\raggedleft
\texttt{wer\_zh\_jieba\_micro}
\end{minipage} & \begin{minipage}[b]{\linewidth}\raggedright
Locale/runtime result
\end{minipage} & \begin{minipage}[b]{\linewidth}\raggedright
Decision
\end{minipage} \\
\midrule\noalign{}
\endhead
\bottomrule\noalign{}
\endlastfoot
Whisper large-v3 & Fixed 15-row contract & 15 & 33.33 & 42.04 & 2
locale-violation rows, 23 simplified chars & Do not promote to 258-row
or selected-300 \\
Whisper large-v3-turbo & Fixed 15-row contract & 15 & 40.36 & 50.87 & 4
locale-violation rows, 48 simplified chars & Retain as bounded
feasibility evidence only; no split-level or selected-300 claim \\
SenseVoiceSmall & Fixed 15-row contract & 15 & 63.12 & 78.98 & 14
locale-violation rows, 209 simplified chars & Reject from full split
until locale policy changes \\
Qwen3-ASR-0.6B & Fixed 15-row contract & 15 & 64.16 & 81.07 & 15
locale-violation rows, 260 simplified chars & Reject from full split
until locale policy changes \\
Qwen3-ASR-1.7B & Bounded load gate & 0 & n/a & n/a & Timeout before
inference at fetch/load after about 60.07s & Retry only after isolated
cache/download plan \\
Gemma 4 E2B/E4B & Local multimodal class probe & 0 & n/a & n/a & Local
Transformers 4.57.6 has no required Gemma 4 multimodal/audio class &
Build isolated prompted multimodal runtime before testing \\
\end{longtable}

Evidence sources:
\texttt{70\_experiments/runs/asr\_candidate\_current\_recheck\_2026\_05\_26/candidate\_current\_recheck\_summary.tsv}
and
\texttt{70\_experiments/runs/asr\_candidate\_current\_recheck\_2026\_05\_26/summary.json}.

Whisper large-v3, Whisper large-v3-turbo, SenseVoiceSmall, and
Qwen3-ASR-0.6B have field-contract evidence but fail strict Taiwan
Traditional Chinese locale promotion criteria. Qwen3-ASR-1.7B remains
stopped before inference at fetch/load, and Gemma 4 E2B/E4B remain
blocked by local multimodal runtime support. The next validation layer
is locale/runtime promotion, not a broader full-split ASR run.

\hypertarget{table-3.-human-reviewed-predictor-table}{%
\subsubsection{Table 3. Human-Reviewed Predictor
Table}\label{table-3.-human-reviewed-predictor-table}}

Table 3 reports aggregate predictor comparison against human-reviewed
\texttt{human\_decision\_change\_yes} labels over 90 reviewed model
assessments. WER/CER are transcript-surface baselines; SRES is the
semantic-risk baseline; CEIS is the proposed decision-stability metric.

Target: \texttt{human\_decision\_change\_yes}, 90 reviewed model
assessments, 16 positive model assessments, clustered within 30 reviewed
audio rows. Diagnostic thresholds are selected on the scoped audit set
for aggregate comparison and are not frozen deployment thresholds.

\begin{longtable}[]{@{}
  >{\raggedright\arraybackslash}p{(\columnwidth - 20\tabcolsep) * \real{0.0732}}
  >{\raggedright\arraybackslash}p{(\columnwidth - 20\tabcolsep) * \real{0.0732}}
  >{\raggedleft\arraybackslash}p{(\columnwidth - 20\tabcolsep) * \real{0.0976}}
  >{\raggedleft\arraybackslash}p{(\columnwidth - 20\tabcolsep) * \real{0.0976}}
  >{\raggedleft\arraybackslash}p{(\columnwidth - 20\tabcolsep) * \real{0.0976}}
  >{\raggedleft\arraybackslash}p{(\columnwidth - 20\tabcolsep) * \real{0.0976}}
  >{\raggedleft\arraybackslash}p{(\columnwidth - 20\tabcolsep) * \real{0.0976}}
  >{\raggedleft\arraybackslash}p{(\columnwidth - 20\tabcolsep) * \real{0.0976}}
  >{\raggedleft\arraybackslash}p{(\columnwidth - 20\tabcolsep) * \real{0.0976}}
  >{\raggedleft\arraybackslash}p{(\columnwidth - 20\tabcolsep) * \real{0.0976}}
  >{\raggedright\arraybackslash}p{(\columnwidth - 20\tabcolsep) * \real{0.0732}}@{}}
\toprule\noalign{}
\begin{minipage}[b]{\linewidth}\raggedright
Predictor
\end{minipage} & \begin{minipage}[b]{\linewidth}\raggedright
Unit
\end{minipage} & \begin{minipage}[b]{\linewidth}\raggedleft
AUC
\end{minipage} & \begin{minipage}[b]{\linewidth}\raggedleft
Row-clustered AUC 95\% CI
\end{minipage} & \begin{minipage}[b]{\linewidth}\raggedleft
Diagnostic threshold
\end{minipage} & \begin{minipage}[b]{\linewidth}\raggedleft
Best F1
\end{minipage} & \begin{minipage}[b]{\linewidth}\raggedleft
Row-clustered F1 95\% CI
\end{minipage} & \begin{minipage}[b]{\linewidth}\raggedleft
Precision
\end{minipage} & \begin{minipage}[b]{\linewidth}\raggedleft
Recall
\end{minipage} & \begin{minipage}[b]{\linewidth}\raggedleft
False negative
\end{minipage} & \begin{minipage}[b]{\linewidth}\raggedright
Paper use
\end{minipage} \\
\midrule\noalign{}
\endhead
\bottomrule\noalign{}
\endlastfoot
WER & 90 model assessments / 30 rows & 0.6964 & 0.5699-0.8125 & 42.59 &
0.4615 & 0.2857-0.6667 & 0.3913 & 0.5625 & 7 & Surface baseline \\
CER & 90 model assessments / 30 rows & 0.7276 & 0.6084-0.8372 & 16.42 &
0.4516 & 0.3158-0.6667 & 0.3043 & 0.8750 & 2 & Surface baseline \\
SRES total & 90 model assessments / 30 rows & 0.8995 & 0.8119-0.9657 &
270.0 & 0.6512 & 0.4667-0.8421 & 0.5185 & 0.8750 & 2 & Semantic-risk
baseline \\
CEIS max & 90 model assessments / 30 rows & 0.9117 & 0.8516-0.9615 & 5.0
& 0.6275 & 0.4681-0.8148 & 0.4571 & 1.0000 & 0 & Proposed
decision-stability metric \\
\end{longtable}

Evidence source:
\texttt{70\_experiments/runs/janus\_300\_high\_stakes\_human\_audit\_selection\_2026\_05\_25/human\_audit\_predictor\_comparison.tsv}.
Row-clustered uncertainty source:
\texttt{70\_experiments/runs/janus\_300\_high\_stakes\_human\_audit\_selection\_2026\_05\_25/human\_audit\_predictor\_clustered\_ci.tsv}.
Leave-one-row-out sensitivity source:
\texttt{70\_experiments/runs/janus\_300\_high\_stakes\_human\_audit\_selection\_2026\_05\_25/human\_audit\_predictor\_leave\_one\_row\_out.tsv}.

CEIS has the strongest point-estimate human-reviewed decision-change AUC
and reaches recall 1.0 at the diagnostic threshold, while SRES achieves
the highest best-threshold F1 and fewer false positives. Row-clustered
AUC intervals overlap for CEIS and SRES, so the evidence supports CEIS
as a conservative decision-stability signal rather than a universally
dominant classifier.

\hypertarget{table-4.-human-reviewed-recovery-policy-table}{%
\subsubsection{Table 4. Human-Reviewed Recovery Policy
Table}\label{table-4.-human-reviewed-recovery-policy-table}}

Table 4 reports five recovery policies evaluated as aggregate replay
under human-reviewed selected-300 labels. The recovery evidence is
aggregate-only and uses the same 30 reviewed rows and 90 reviewed model
assessments as the predictor table.

Evidence mode: human-reviewed selected-300, 30 reviewed rows and 90
reviewed model assessments.

\begin{longtable}[]{@{}
  >{\raggedright\arraybackslash}p{(\columnwidth - 20\tabcolsep) * \real{0.0714}}
  >{\raggedleft\arraybackslash}p{(\columnwidth - 20\tabcolsep) * \real{0.0952}}
  >{\raggedleft\arraybackslash}p{(\columnwidth - 20\tabcolsep) * \real{0.0952}}
  >{\raggedleft\arraybackslash}p{(\columnwidth - 20\tabcolsep) * \real{0.0952}}
  >{\raggedleft\arraybackslash}p{(\columnwidth - 20\tabcolsep) * \real{0.0952}}
  >{\raggedleft\arraybackslash}p{(\columnwidth - 20\tabcolsep) * \real{0.0952}}
  >{\raggedleft\arraybackslash}p{(\columnwidth - 20\tabcolsep) * \real{0.0952}}
  >{\raggedleft\arraybackslash}p{(\columnwidth - 20\tabcolsep) * \real{0.0952}}
  >{\raggedleft\arraybackslash}p{(\columnwidth - 20\tabcolsep) * \real{0.0952}}
  >{\raggedleft\arraybackslash}p{(\columnwidth - 20\tabcolsep) * \real{0.0952}}
  >{\raggedright\arraybackslash}p{(\columnwidth - 20\tabcolsep) * \real{0.0714}}@{}}
\toprule\noalign{}
\begin{minipage}[b]{\linewidth}\raggedright
Policy
\end{minipage} & \begin{minipage}[b]{\linewidth}\raggedleft
Model assessments
\end{minipage} & \begin{minipage}[b]{\linewidth}\raggedleft
Unsafe downrouting
\end{minipage} & \begin{minipage}[b]{\linewidth}\raggedleft
High-risk missed
\end{minipage} & \begin{minipage}[b]{\linewidth}\raggedleft
Critical miss
\end{minipage} & \begin{minipage}[b]{\linewidth}\raggedleft
Recovery budget
\end{minipage} & \begin{minipage}[b]{\linewidth}\raggedleft
Row-clustered budget 95\% CI
\end{minipage} & \begin{minipage}[b]{\linewidth}\raggedleft
Severe misses eliminated
\end{minipage} & \begin{minipage}[b]{\linewidth}\raggedleft
Row-clustered severe-miss 95\% CI
\end{minipage} & \begin{minipage}[b]{\linewidth}\raggedleft
Triggers per severe miss eliminated
\end{minipage} & \begin{minipage}[b]{\linewidth}\raggedright
Scoped interpretation
\end{minipage} \\
\midrule\noalign{}
\endhead
\bottomrule\noalign{}
\endlastfoot
No recovery & 90 & 29 & 6 & 1 & 0.0000 & 0.0000-0.0000 & 0 & 2-13
remaining severe misses & n/a & Baseline decision risk \\
Calibrated-confidence unavailable & 90 & 29 & 6 & 1 & 0.0000 &
0.0000-0.0000 & 0 & 2-13 remaining severe misses & n/a & No calibrated
confidence trigger available \\
SRES-triggered recovery & 90 & 24 & 0 & 0 & 0.3889 & 0.3000-0.4889 & 7 &
0-0 remaining severe misses & 5.0 & Semantic-risk recovery baseline \\
CEIS-triggered conservative action & 90 & 24 & 0 & 0 & 0.3889 &
0.3000-0.4889 & 7 & 0-0 remaining severe misses & 5.0 & Matches SRES
severe-miss elimination at same budget in replay \\
CEIS ensemble arbitration & 90 & 24 & 0 & 0 & 0.5222 & 0.4000-0.6556 & 7
& 0-0 remaining severe misses & 6.7 & Adds abstention/interval behavior
at higher budget \\
\end{longtable}

Evidence sources:
\texttt{70\_experiments/runs/janus\_300\_high\_stakes\_recovery\_human\_reviewed\_2026\_05\_26/policy\_comparison.tsv}
and
\texttt{70\_experiments/runs/janus\_300\_high\_stakes\_recovery\_human\_reviewed\_2026\_05\_26/summary.json}.
Row-clustered uncertainty source:
\texttt{70\_experiments/runs/janus\_300\_high\_stakes\_recovery\_human\_reviewed\_2026\_05\_26/policy\_comparison\_clustered\_ci.tsv}.
Leave-one-row-out sensitivity source:
\texttt{70\_experiments/runs/janus\_300\_high\_stakes\_recovery\_human\_reviewed\_2026\_05\_26/policy\_comparison\_leave\_one\_row\_out.tsv}.
Fixed-budget frontier source:
\texttt{70\_experiments/runs/janus\_300\_high\_stakes\_recovery\_human\_reviewed\_2026\_05\_26/fixed\_budget\_recovery\_frontier.tsv}.

Without recovery, the reviewed evidence contains 6 high-risk misses and
1 critical miss. SRES-triggered recovery and CEIS-triggered conservative
action both eliminate high-risk missed and critical miss counts in
aggregate policy replay at the same 0.3889 trigger budget, corresponding
to 35 triggers for 7 severe missed outcomes eliminated. CEIS ensemble
arbitration preserves this 0/0 result while introducing abstention
behavior at a higher 0.5222 budget. The policy replay eliminates the
most severe missed-risk outcomes under the scoped labels, while residual
unsafe downrouting remains at 24 and requires separate governance.

A fixed-budget replay provides an additional operating view. At 10\%,
20\%, 30\%, and 40\% requested trigger budgets, CEIS-ranked conservative
replay leaves 0 high-risk misses and 0 critical misses under the scoped
labels; the 10\% point uses 9 triggers. SRES-ranked conservative replay
leaves 4, 2, 2, and 0 severe missed outcomes at the same requested
budgets, reaching 0 severe misses only when all 35 eligible SRES
triggers are used. This frontier supports CEIS as a conservative
decision-stability signal while preserving the Table 4 result that the
diagnostic SRES and CEIS policies tie at the selected 0.3889 budget.
Appendix Table A1 reports the fixed-budget replay explicitly as an
operating-point frontier rather than a deployment-threshold claim.

\hypertarget{figure-package}{%
\subsection{Figure Package}\label{figure-package}}

The figure package is generated from aggregate-only inputs by
\texttt{80\_semantic\_risk\_asr/paper/generate\_paper\_figures.py}. The
generated SVGs and PDF exports and their privacy boundary are listed in
\texttt{80\_semantic\_risk\_asr/paper/figures/}. The aggregate artifact
manifest is generated by
\texttt{80\_semantic\_risk\_asr/paper/build\_artifact\_manifest.py} and
written to
\texttt{80\_semantic\_risk\_asr/paper/artifact\_manifest.tsv}.

\begin{longtable}[]{@{}
  >{\raggedright\arraybackslash}p{(\columnwidth - 8\tabcolsep) * \real{0.2000}}
  >{\raggedright\arraybackslash}p{(\columnwidth - 8\tabcolsep) * \real{0.2000}}
  >{\raggedright\arraybackslash}p{(\columnwidth - 8\tabcolsep) * \real{0.2000}}
  >{\raggedright\arraybackslash}p{(\columnwidth - 8\tabcolsep) * \real{0.2000}}
  >{\raggedright\arraybackslash}p{(\columnwidth - 8\tabcolsep) * \real{0.2000}}@{}}
\toprule\noalign{}
\begin{minipage}[b]{\linewidth}\raggedright
Figure
\end{minipage} & \begin{minipage}[b]{\linewidth}\raggedright
File
\end{minipage} & \begin{minipage}[b]{\linewidth}\raggedright
Caption
\end{minipage} & \begin{minipage}[b]{\linewidth}\raggedright
Source
\end{minipage} & \begin{minipage}[b]{\linewidth}\raggedright
Privacy boundary
\end{minipage} \\
\midrule\noalign{}
\endhead
\bottomrule\noalign{}
\endlastfoot
F1. CDS-ASR pipeline & \texttt{figures/f1\_cds\_asr\_pipeline.svg} &
CDS-ASR converts audio into ASR hypotheses and runtime signals, extracts
risk atoms, generates plausible variants, scores SRES/CEIS, and applies
constrained recovery or conservative action. & Method section & No
transcript text or row content \\
F2. Evidence boundary & \texttt{figures/f2\_evidence\_boundary.svg} &
The manuscript separates 258-row split/model-comparison evidence,
selected-300 provenance evidence, and selected-300 human-reviewed
predictor/recovery evidence. & publishable evidence summary & Aggregate
status only \\
F3. Predictor AUC & \texttt{figures/f3\_predictor\_auc.svg} & CEIS has
the highest human-reviewed decision-change AUC in the scoped
selected-300 audit, while SRES remains strongest on best-threshold F1;
Table 3 reports row-clustered uncertainty. &
\texttt{human\_audit\_predictor\_comparison.tsv} & Aggregate predictor
metrics \\
F4. Recovery outcomes & \texttt{figures/f4\_recovery\_outcomes.svg} &
SRES-triggered recovery and CEIS-triggered conservative action both
eliminate high-risk missed and critical miss counts in aggregate replay
at the same 0.3889 budget. & \texttt{policy\_comparison.tsv} & Aggregate
policy counts \\
F5. Model lane state & \texttt{figures/f5\_model\_lane\_state.svg} &
Main comparable split evidence, locale-gated candidates, and
runtime-blocked probes remain separate until promotion gates are
satisfied. & main/candidate aggregate summaries & Aggregate lane
state \\
F6. N-ladder & \texttt{figures/f6\_n\_ladder.svg} & The evidence chain
separates 258 rows, selected-300 provenance, 30 reviewed rows, and 90
clustered model assessments. & method evidence units & Aggregate counts
only \\
F7. Budget-risk frontier &
\texttt{figures/f7\_budget\_risk\_frontier.svg} & Fixed-budget policy
replay shows the trigger-budget tradeoff needed to eliminate severe
missed outcomes under scoped labels. &
\texttt{fixed\_budget\_recovery\_frontier.tsv} & Aggregate policy
counts \\
\end{longtable}

Submission-package PDF exports are generated with the same basename
under \texttt{80\_semantic\_risk\_asr/paper/figures/}.

\hypertarget{discussion}{%
\subsection{Discussion}\label{discussion}}

The evidence supports a consequence-centered ASR evaluation claim:
correct WER/CER reporting is necessary, but transcript similarity alone
does not capture high-stakes decision instability. The reviewed
selected-300 predictor table shows that SRES and CEIS align more
strongly with human-reviewed decision-change labels than WER/CER in this
scoped evidence layer.

The model comparison should be presented as an evidence layer, not as a
model leaderboard. The partial encoder is the strongest current ASR
hypothesis generator on the comparable split, and the selected-300
human-reviewed evidence shows how downstream decision evidence refines
transcript-centered model comparison.

The recovery result supports a scoped replay claim. Under human-reviewed
selected-300 labels, risk-triggered conservative policies, including
CEIS-triggered conservative action, eliminate high-risk missed and
critical miss counts in aggregate replay. SRES-triggered recovery and
CEIS-triggered conservative action reach this result at the same 0.3889
trigger budget, while ensemble arbitration adds interval/abstention
behavior as a higher-budget governance option. CEIS is therefore best
presented as the more conservative decision-instability signal at the
predictor layer, while SRES remains a strong semantic-risk recovery
baseline.

Improving ASR remains necessary, but stronger transcript models do not
replace decision-stability analysis. The paper studies residual
instability after transcript scoring: whether plausible ASR alternatives
around decision atoms would change escalation, routing, or conservative
action. This makes CDS-ASR a safety layer for the remaining decision
interval, not a substitute for ASR model improvement.

The aggregate-only artifact boundary is also a contribution to
reproducible governance. The paper does not claim full public row-level
reproducibility. Instead, it provides aggregate reproducibility,
operation records, validation gates, and consistency checks under a
privacy-preserving audit boundary.

The candidate lane is already bounded by evidence. Whisper large-v3,
Whisper large-v3-turbo, SenseVoiceSmall, and Qwen3-ASR-0.6B have
small-gate evidence but fail strict Taiwan Traditional Chinese locale
promotion criteria. Qwen3-ASR-1.7B is stopped by fetch/load timeout, and
Gemma 4 E2B/E4B require an isolated multimodal runtime. The next
validation step is locale/runtime validation, not a broader full-split
ASR run.

\hypertarget{claim-registry}{%
\subsection{Claim Registry}\label{claim-registry}}

\begin{longtable}[]{@{}
  >{\raggedright\arraybackslash}p{(\columnwidth - 8\tabcolsep) * \real{0.2000}}
  >{\raggedright\arraybackslash}p{(\columnwidth - 8\tabcolsep) * \real{0.2000}}
  >{\raggedright\arraybackslash}p{(\columnwidth - 8\tabcolsep) * \real{0.2000}}
  >{\raggedright\arraybackslash}p{(\columnwidth - 8\tabcolsep) * \real{0.2000}}
  >{\raggedright\arraybackslash}p{(\columnwidth - 8\tabcolsep) * \real{0.2000}}@{}}
\toprule\noalign{}
\begin{minipage}[b]{\linewidth}\raggedright
Claim
\end{minipage} & \begin{minipage}[b]{\linewidth}\raggedright
Scope
\end{minipage} & \begin{minipage}[b]{\linewidth}\raggedright
Evidence artifact
\end{minipage} & \begin{minipage}[b]{\linewidth}\raggedright
Statistic
\end{minipage} & \begin{minipage}[b]{\linewidth}\raggedright
Limitation / defense
\end{minipage} \\
\midrule\noalign{}
\endhead
\bottomrule\noalign{}
\endlastfoot
CEIS has the strongest AUC among reported predictors & 90 clustered
model assessments from 30 reviewed rows &
\texttt{human\_audit\_predictor\_comparison.tsv} & AUC 0.9117 &
Row-clustered CI needed; selected high-stakes audit surface \\
CEIS reaches a zero-FN operating point & scoped selected-300 diagnostic
threshold & \texttt{human\_audit\_predictor\_comparison.tsv} & FN 0,
recall 1.0000 & Retrospective diagnostic threshold, not frozen
deployment threshold \\
SRES and CEIS conservative policies eliminate severe misses in replay &
aggregate policy replay over reviewed assessments &
\texttt{policy\_comparison.tsv} & high-risk missed 0, critical miss 0 &
Replay evidence, not live causal deployment; policies tie at 0.3889
budget \\
Partial encoder is the strongest current ASR hypothesis generator &
258-row split/model-comparison layer &
\texttt{asr\_cds\_proxy\_comparison.tsv} & \texttt{cer\_zh\_micro}
15.04, unsafe downrouting 7, high-risk missed 4 & Split-level
model-comparison context only \\
Selected-300 is a high-stakes audit surface & enriched selection
provenance & \texttt{human\_audit\_selection\_summary.json},
\texttt{selection\_strata.tsv} & 300 candidates, 30 reviewed rows, 90
assessments & Not prevalence-preserving; selection enrichment
disclosed \\
Aggregate-only release supports reviewer-visible auditability & release
and operation-record layer & validation summaries, evidence matrices,
operation records, figure scripts & gate state: roadmap complete,
publishable ready, consistency 26/26 & No public row-level transcript
reproducibility \\
\end{longtable}

\hypertarget{ethics-privacy-and-intended-use}{%
\subsection{Ethics, Privacy, And Intended
Use}\label{ethics-privacy-and-intended-use}}

This study treats transcript-bearing speech data as sensitive
operational content. NIH human-subjects guidance treats the use, study,
analysis, or generation of identifiable private information as a
human-subjects trigger under the cited federal definition (National
Institutes of Health 2024). This paper therefore keeps raw audio,
transcript-bearing hypotheses, reviewer notes, row identifiers, local
response sheets, and runtime logs outside the release boundary. Before
any external data release or deployment claim, the institutional review,
exemption, data-use, retention, encryption, and access control status
must be documented.

CDS-ASR is intended for ASR safety audit, risk-aware routing,
manual-review prioritization, conservative escalation, and machine
abstention. It is not intended for automatic guilt or fraud
determination, automatic account freezing, service denial, punitive
action, or direct law-enforcement reporting without human review.
``Conservative machine action'' means preserving uncertainty, raising
review priority, abstaining, or requiring human confirmation.

This intended-use boundary aligns the paper with risk-management
governance rather than autonomous adverse decision-making. NIST
describes the AI RMF as a voluntary framework for managing AI risks to
individuals, organizations, and society and for incorporating
trustworthiness considerations into AI design, development, use, and
evaluation (National Institute of Standards and Technology 2026). For
cross-border or future deployment contexts, the EU AI Act illustrates
the broader regulatory direction toward risk-based obligations and
high-risk AI governance (European Commission 2026). This manuscript does
not provide legal compliance analysis; it provides a scoped
consequence-centered audit method and a privacy-preserving release
boundary.

\hypertarget{limitations-and-threats-to-validity}{%
\subsection{Limitations And Threats To
Validity}\label{limitations-and-threats-to-validity}}

The human-reviewed evidence is scoped to 30 rows and 90 model
assessments. It supports a focused high-stakes audit claim, not a
population-level deployment claim.

The 90 model assessments are clustered within 30 audio rows and should
not be treated as 90 independent calls. The current submission-prep
package reports row-clustered bootstrap intervals and leave-one-row-out
sensitivity tables for predictor and recovery metrics; these remain
uncertainty descriptions for a scoped audit rather than population-level
deployment intervals.

The number of positive decision-change cases is limited. Table 3 has 16
positive model assessments, so AUC and threshold behavior should be
interpreted with uncertainty.

The selected-300 audit surface is deliberately enriched for high-stakes
signals. It supports decision-stability behavior within the selected
audit boundary, not population prevalence of ASR-induced harm across all
calls.

The reported best thresholds are diagnostic operating points on the
scoped reviewed audit. They can overstate deployment performance unless
threshold selection is later frozen on a separate development set.

Human labels come from a completed expert audit rather than
multi-annotator adjudication. No second-reviewer blinded spot-check is
included in the current submission-prep package, so inter-annotator
agreement is not claimed. This is a single-expert audit boundary, not a
multi-annotator reliability claim.

Recovery evidence is aggregate-only. This protects transcript-bearing
call and review content, but limits external row-level reproducibility.

CEIS depends on generated plausible variants and risk atom weights.
Missed variants or misweighted atoms can affect instability scoring. The
current aggregate package records risk-atom coverage and a
variant-coverage status table, but a full source-specific aggregate
variant generation audit remains a planned validation extension.

The Taiwan Traditional Chinese locale policy is strict by design.
Candidate model rejection may reflect deployment-locale mismatch, not
universal ASR inferiority.

Recovery policies eliminate high-risk missed and critical miss counts in
aggregate replay, but Table 4 still leaves unsafe downrouting count at
24 after SRES/CEIS conservative recovery. The evidence supports scoped
severe-miss reduction, not elimination of all safety risk.

The study does not include a deployment trial and does not estimate
population prevalence of ASR-induced harm across all calls. It evaluates
whether decision-critical ASR risk can be detected and mitigated within
a deliberately enriched high-stakes audit surface.

\hypertarget{appendix-artifact-availability}{%
\subsection{Appendix / Artifact
Availability}\label{appendix-artifact-availability}}

\hypertarget{artifact-availability-statement}{%
\subsubsection{Artifact Availability
Statement}\label{artifact-availability-statement}}

The repository release includes aggregate run records, validation
summaries, metric tables, evidence matrices, figure-generation code, and
paper-facing documentation. Raw audio, raw transcripts, selected row
identifiers, audio identifiers, model hypotheses, transcript-bearing
runtime logs, reviewer response sheets, reviewer notes, and model
weights are not released because they may contain sensitive call or
review content.

Because transcript-bearing materials may contain sensitive call content,
reproducibility is provided through aggregate validation artifacts,
operation records, manifest checks, and consistency audits rather than
through raw row release. This provides reviewer-visible auditability for
model comparison, metric policy, selected-300 human-review status,
predictor analysis, recovery-policy evaluation, candidate-lane
boundaries, and consistency checks while preserving the local-only
boundary for sensitive materials.

The reviewer package should include an aggregate manifest with artifact
path, role, privacy class, generator, source inputs, SHA256, source git
commit, timestamp, Python version, metric-library versions, tokenizer
policy, locale normalization policy, decoding parameters where
available, random seeds where used, and hardware summary. The manifest
records auditability metadata without adding transcript-bearing row
content.

\hypertarget{reviewer-reproducible-aggregate-artifacts}{%
\subsubsection{Reviewer-Reproducible Aggregate
Artifacts}\label{reviewer-reproducible-aggregate-artifacts}}

\begin{itemize}
\tightlist
\item
  Artifact manifest:
  \texttt{80\_semantic\_risk\_asr/paper/artifact\_manifest.tsv}
\item
  Artifact manifest generator:
  \texttt{80\_semantic\_risk\_asr/paper/build\_artifact\_manifest.py}
\item
  Experiment registry: \texttt{70\_experiments/registry.tsv}
\item
  Main 258-row comparison:
  \texttt{70\_experiments/runs/janus\_258\_test\_split\_asr\_cds\_proxy/asr\_cds\_proxy\_comparison.tsv}
\item
  Metric-definition audit:
  \texttt{70\_experiments/runs/wer\_metric\_audit\_2026\_05\_25/journal\_compliance\_summary.json}
\item
  Selected-300 human audit refresh:
  \texttt{70\_experiments/runs/janus\_300\_high\_stakes\_human\_audit\_selection\_2026\_05\_25/human\_audit\_refresh\_summary.json}
\item
  Human-reviewed predictor evidence:
  \texttt{70\_experiments/runs/janus\_300\_high\_stakes\_human\_audit\_selection\_2026\_05\_25/human\_audit\_predictor\_comparison.tsv}
\item
  Human-reviewed predictor row-clustered CI:
  \texttt{70\_experiments/runs/janus\_300\_high\_stakes\_human\_audit\_selection\_2026\_05\_25/human\_audit\_predictor\_clustered\_ci.tsv}
\item
  Human-reviewed predictor leave-one-row-out sensitivity:
  \texttt{70\_experiments/runs/janus\_300\_high\_stakes\_human\_audit\_selection\_2026\_05\_25/human\_audit\_predictor\_leave\_one\_row\_out.tsv}
\item
  Human-reviewed recovery evidence:
  \texttt{70\_experiments/runs/janus\_300\_high\_stakes\_recovery\_human\_reviewed\_2026\_05\_26/policy\_comparison.tsv}
\item
  Human-reviewed recovery row-clustered CI:
  \texttt{70\_experiments/runs/janus\_300\_high\_stakes\_recovery\_human\_reviewed\_2026\_05\_26/policy\_comparison\_clustered\_ci.tsv}
\item
  Human-reviewed recovery leave-one-row-out sensitivity:
  \texttt{70\_experiments/runs/janus\_300\_high\_stakes\_recovery\_human\_reviewed\_2026\_05\_26/policy\_comparison\_leave\_one\_row\_out.tsv}
\item
  Human-reviewed recovery fixed-budget frontier:
  \texttt{70\_experiments/runs/janus\_300\_high\_stakes\_recovery\_human\_reviewed\_2026\_05\_26/fixed\_budget\_recovery\_frontier.tsv}
\item
  Claim registry:
  \texttt{70\_experiments/runs/postdoc\_evidence\_chain\_2026\_05\_25/claim\_registry.tsv}
\item
  CEIS method summary:
  \texttt{70\_experiments/runs/postdoc\_evidence\_chain\_2026\_05\_25/ceis\_method\_summary.tsv}
\item
  Selection provenance summary:
  \texttt{70\_experiments/runs/janus\_300\_high\_stakes\_human\_audit\_selection\_2026\_05\_25/selection\_provenance\_summary.tsv}
\item
  Counterfactual variant coverage status:
  \texttt{70\_experiments/runs/janus\_300\_high\_stakes\_human\_audit\_selection\_2026\_05\_25/counterfactual\_variant\_coverage\_summary.tsv}
\item
  Post-review evidence checklist:
  \texttt{70\_experiments/runs/janus\_300\_high\_stakes\_human\_audit\_selection\_2026\_05\_25/human\_audit\_post\_review\_evidence\_summary.json}
\item
  Publishable evidence audit:
  \texttt{70\_experiments/runs/postdoc\_evidence\_chain\_2026\_05\_25/publishable\_evidence\_completion\_summary.json}
\item
  Consequence evidence matrix:
  \texttt{70\_experiments/runs/postdoc\_evidence\_chain\_2026\_05\_25/consequence\_evidence\_matrix.tsv}
\item
  Roadmap completion audit:
  \texttt{70\_experiments/runs/postdoc\_evidence\_chain\_2026\_05\_25/postdoc\_roadmap\_completion\_summary.json}
\item
  Evidence-chain consistency audit:
  \texttt{70\_experiments/runs/postdoc\_evidence\_chain\_2026\_05\_25/evidence\_chain\_consistency\_summary.json}
\item
  Candidate bounded recheck:
  \texttt{70\_experiments/runs/asr\_candidate\_current\_recheck\_2026\_05\_26/summary.json}
\end{itemize}

\hypertarget{appendix-table-a1.-fixed-budget-recovery-frontier}{%
\subsubsection{Appendix Table A1. Fixed-Budget Recovery
Frontier}\label{appendix-table-a1.-fixed-budget-recovery-frontier}}

The fixed-budget frontier is a retrospective aggregate replay over the
same 90 reviewed model assessments. It reports what remains after
applying ranked conservative triggers at requested trigger budgets.
Requested 40\% budget maps to 35 eligible triggers, so the observed
budget is 0.3889.

\begin{longtable}[]{@{}
  >{\raggedright\arraybackslash}p{(\columnwidth - 18\tabcolsep) * \real{0.0769}}
  >{\raggedleft\arraybackslash}p{(\columnwidth - 18\tabcolsep) * \real{0.1026}}
  >{\raggedleft\arraybackslash}p{(\columnwidth - 18\tabcolsep) * \real{0.1026}}
  >{\raggedleft\arraybackslash}p{(\columnwidth - 18\tabcolsep) * \real{0.1026}}
  >{\raggedleft\arraybackslash}p{(\columnwidth - 18\tabcolsep) * \real{0.1026}}
  >{\raggedleft\arraybackslash}p{(\columnwidth - 18\tabcolsep) * \real{0.1026}}
  >{\raggedleft\arraybackslash}p{(\columnwidth - 18\tabcolsep) * \real{0.1026}}
  >{\raggedleft\arraybackslash}p{(\columnwidth - 18\tabcolsep) * \real{0.1026}}
  >{\raggedleft\arraybackslash}p{(\columnwidth - 18\tabcolsep) * \real{0.1026}}
  >{\raggedleft\arraybackslash}p{(\columnwidth - 18\tabcolsep) * \real{0.1026}}@{}}
\toprule\noalign{}
\begin{minipage}[b]{\linewidth}\raggedright
Score metric
\end{minipage} & \begin{minipage}[b]{\linewidth}\raggedleft
Requested budget
\end{minipage} & \begin{minipage}[b]{\linewidth}\raggedleft
Triggered
\end{minipage} & \begin{minipage}[b]{\linewidth}\raggedleft
Observed budget
\end{minipage} & \begin{minipage}[b]{\linewidth}\raggedleft
Unsafe downrouting
\end{minipage} & \begin{minipage}[b]{\linewidth}\raggedleft
High-risk missed
\end{minipage} & \begin{minipage}[b]{\linewidth}\raggedleft
Critical miss
\end{minipage} & \begin{minipage}[b]{\linewidth}\raggedleft
Severe missed
\end{minipage} & \begin{minipage}[b]{\linewidth}\raggedleft
Severe misses eliminated
\end{minipage} & \begin{minipage}[b]{\linewidth}\raggedleft
Triggers per severe miss eliminated
\end{minipage} \\
\midrule\noalign{}
\endhead
\bottomrule\noalign{}
\endlastfoot
SRES total & 0.1000 & 9 & 0.1000 & 28 & 4 & 0 & 4 & 3 & 3.0 \\
SRES total & 0.2000 & 18 & 0.2000 & 26 & 2 & 0 & 2 & 5 & 3.6 \\
SRES total & 0.3000 & 27 & 0.3000 & 26 & 2 & 0 & 2 & 5 & 5.4 \\
SRES total & 0.4000 & 35 & 0.3889 & 24 & 0 & 0 & 0 & 7 & 5.0 \\
CEIS & 0.1000 & 9 & 0.1000 & 24 & 0 & 0 & 0 & 7 & 1.2857 \\
CEIS & 0.2000 & 18 & 0.2000 & 24 & 0 & 0 & 0 & 7 & 2.5714 \\
CEIS & 0.3000 & 27 & 0.3000 & 24 & 0 & 0 & 0 & 7 & 3.8571 \\
CEIS & 0.4000 & 35 & 0.3889 & 24 & 0 & 0 & 0 & 7 & 5.0 \\
\end{longtable}

Evidence source:
\texttt{70\_experiments/runs/janus\_300\_high\_stakes\_recovery\_human\_reviewed\_2026\_05\_26/fixed\_budget\_recovery\_frontier.tsv}.

\hypertarget{operation-records}{%
\subsubsection{Operation Records}\label{operation-records}}

The selected-300 reviewer route is documented through aggregate-only
operation records, including response closeout, response apply logs,
reviewer work order, post-review sequence, operation-record audit, and
consistency audit. These records provide command-level reproducibility
without exposing transcript or row content.

Key records:

\begin{itemize}
\tightlist
\item
  \texttt{70\_experiments/runs/janus\_300\_high\_stakes\_human\_audit\_selection\_2026\_05\_25/human\_audit\_response\_closeout\_summary.json}
\item
  \texttt{70\_experiments/runs/janus\_300\_high\_stakes\_human\_audit\_selection\_2026\_05\_25/human\_audit\_batch\_response\_apply\_log.tsv}
\item
  \texttt{70\_experiments/runs/janus\_300\_high\_stakes\_human\_audit\_selection\_2026\_05\_25/human\_audit\_review\_work\_order\_summary.json}
\item
  \texttt{70\_experiments/runs/janus\_300\_high\_stakes\_human\_audit\_selection\_2026\_05\_25/human\_audit\_post\_review\_sequence\_summary.json}
\item
  \texttt{70\_experiments/runs/janus\_300\_high\_stakes\_human\_audit\_selection\_2026\_05\_25/human\_audit\_operation\_record\_summary.json}
\end{itemize}

\hypertarget{privacy-and-local-only-boundary}{%
\subsubsection{Privacy And Local-Only
Boundary}\label{privacy-and-local-only-boundary}}

The following stay local or ignored:

\begin{itemize}
\tightlist
\item
  raw audio;
\item
  raw transcripts;
\item
  selected sample IDs and audio IDs;
\item
  transcript-bearing audit sheets and response sheets;
\item
  reviewer notes;
\item
  raw model hypotheses and predictions;
\item
  transcript-bearing runtime logs;
\item
  checkpoints, adapters, and model weights;
\item
  local downloaded reviewer packets and zip files.
\end{itemize}

The tracked repo preserves aggregate counts, run records, validation
summaries, metric tables, and paper-facing evidence matrices. This
boundary lets reviewers audit the evidence chain while protecting
transcript-bearing material.

\hypertarget{validation-gate-commands}{%
\subsection{Validation Gate Commands}\label{validation-gate-commands}}

Run these after manuscript packaging changes:

\begin{Shaded}
\begin{Highlighting}[]
\ExtensionTok{.venv/bin/python}\NormalTok{ 80\_semantic\_risk\_asr/scoring/audit\_postdoc\_roadmap\_completion.py }\AttributeTok{{-}{-}output{-}json}\NormalTok{ /tmp/cib\_roadmap\_completion\_check.json }\AttributeTok{{-}{-}output{-}tsv}\NormalTok{ /tmp/cib\_roadmap\_completion\_check.tsv}
\ExtensionTok{.venv/bin/python}\NormalTok{ 80\_semantic\_risk\_asr/scoring/audit\_publishable\_evidence\_chain.py }\AttributeTok{{-}{-}output{-}json}\NormalTok{ /tmp/cib\_publishable\_check.json }\AttributeTok{{-}{-}output{-}tsv}\NormalTok{ /tmp/cib\_publishable\_check.tsv}
\ExtensionTok{.venv/bin/python}\NormalTok{ 80\_semantic\_risk\_asr/scoring/audit\_evidence\_chain\_consistency.py }\AttributeTok{{-}{-}output{-}json}\NormalTok{ /tmp/cib\_consistency\_check.json }\AttributeTok{{-}{-}output{-}tsv}\NormalTok{ /tmp/cib\_consistency\_check.tsv}
\end{Highlighting}
\end{Shaded}

Expected current gate state:

\begin{itemize}
\tightlist
\item
  roadmap completion: \texttt{roadmap\_complete=true},
  \texttt{blocking\_gate=none};
\item
  publishable evidence: \texttt{publishable\_ready=true};
\item
  consistency audit: \texttt{26/26} checks pass,
  \texttt{failed\_checks={[}{]}}.
\end{itemize}

\hypertarget{scope-control-for-additional-experiments}{%
\subsection{Scope Control For Additional
Experiments}\label{scope-control-for-additional-experiments}}

Do not run new full-split ASR experiments now. Candidate models can move
to 258-row or selected-300 only after strict Taiwan Traditional Chinese
locale policy is satisfied or an isolated Gemma 4 multimodal runtime
exists. Until then, the active work is manuscript drafting,
results-table packaging, artifact availability, citation completion, and
submission readiness.

\hypertarget{refs}{}
\begin{CSLReferences}{1}{0}
\leavevmode\vadjust pre{\hypertarget{ref-aws_contact_lens_analytics_2026}{}}%
Amazon Web Services. 2026a. {``Analyze Conversations Using
Conversational Analytics in Connect Customer Contact Lens.''}
\url{https://docs.aws.amazon.com/connect/latest/adminguide/analyze-conversations.html}.

\leavevmode\vadjust pre{\hypertarget{ref-aws_connect_contact_lens_2026}{}}%
---------. 2026b. {``Contact Center Conversational Analytics - Amazon
Connect.''}
\url{https://aws.amazon.com/products/connect/customer/conversational-analytics/}.

\leavevmode\vadjust pre{\hypertarget{ref-angelopoulos2021conformal}{}}%
Angelopoulos, Anastasios N., and Stephen Bates. 2021. {``A Gentle
Introduction to Conformal Prediction and Distribution-Free Uncertainty
Quantification.''} \url{https://arxiv.org/abs/2107.07511}.

\leavevmode\vadjust pre{\hypertarget{ref-chow1970reject}{}}%
Chow, C. K. 1970. {``On Optimum Recognition Error and Reject
Tradeoff.''} \emph{IEEE Transactions on Information Theory} 16 (1):
41--46. \url{https://doi.org/10.1109/TIT.1970.1054406}.

\leavevmode\vadjust pre{\hypertarget{ref-eu_ai_act_2026}{}}%
European Commission. 2026. {``AI Act.''}
\url{https://digital-strategy.ec.europa.eu/en/policies/regulatory-framework-ai}.

\leavevmode\vadjust pre{\hypertarget{ref-fbi_ic3_2025_report_2026}{}}%
Federal Bureau of Investigation. 2026a. {``2025 Internet Crime
Report.''} \url{https://www.fbi.gov/file-repository/2025_ic3report.pdf}.

\leavevmode\vadjust pre{\hypertarget{ref-fbi_crypto_ai_scams_2026}{}}%
---------. 2026b. {``Cryptocurrency and AI Scams Bilk Americans of
Billions.''}
\url{https://www.fbi.gov/news/press-releases/cryptocurrency-and-ai-scams-bilk-americans-of-billions}.

\leavevmode\vadjust pre{\hypertarget{ref-geifman2017selective}{}}%
Geifman, Yonatan, and Ran El-Yaniv. 2017. {``Selective Classification
for Deep Neural Networks.''} In \emph{Advances in Neural Information
Processing Systems 30}.
\url{https://papers.neurips.cc/paper/7073-selective-classification-for-deep-neural-networks}.

\leavevmode\vadjust pre{\hypertarget{ref-geifman2019selectivenet}{}}%
---------. 2019. {``SelectiveNet: A Deep Neural Network with an
Integrated Reject Option.''} In \emph{Proceedings of the 36th
International Conference on Machine Learning}, 97:2151--59. Proceedings
of Machine Learning Research. PMLR.
\url{https://proceedings.mlr.press/v97/geifman19a.html}.

\leavevmode\vadjust pre{\hypertarget{ref-kim2021semanticdistance}{}}%
Kim, Suyoun, Abhinav Arora, Duc Le, Ching-Feng Yeh, Christian Fuegen,
Ozlem Kalinli, and Michael L. Seltzer. 2021. {``Semantic Distance: A New
Metric for ASR Performance Analysis Towards Spoken Language
Understanding.''} In \emph{Interspeech 2021}, 1977--81.
\url{https://doi.org/10.21437/Interspeech.2021-1929}.

\leavevmode\vadjust pre{\hypertarget{ref-miner2020psychotherapy_asr}{}}%
Miner, Adam S., Albert Haque, Jason A. Fries, et al. 2020. {``Assessing
the Accuracy of Automatic Speech Recognition for Psychotherapy.''}
\emph{Npj Digital Medicine} 3: 82.
\url{https://doi.org/10.1038/s41746-020-0285-8}.

\leavevmode\vadjust pre{\hypertarget{ref-naderi2024llmconfidence}{}}%
Naderi, Maryam, Enno Hermann, Alexandre Nanchen, Sevada Hovsepyan, and
Mathew Magimai.-Doss. 2024. {``Towards Interfacing Large Language Models
with ASR Systems Using Confidence Measures and Prompting.''} In
\emph{Interspeech 2024}, 2980--84.
\url{https://doi.org/10.21437/Interspeech.2024-989}.

\leavevmode\vadjust pre{\hypertarget{ref-nist_ai_rmf_2026}{}}%
National Institute of Standards and Technology. 2026. {``AI Risk
Management Framework.''}
\url{https://www.nist.gov/itl/ai-risk-management-framework}.

\leavevmode\vadjust pre{\hypertarget{ref-nih_human_subjects_research_2024}{}}%
National Institutes of Health. 2024. {``Definition of Human Subjects
Research.''}
\url{https://grants.nih.gov/policy-and-compliance/policy-topics/human-subjects/research}.

\leavevmode\vadjust pre{\hypertarget{ref-npa_165_antifraud_hotline_2024}{}}%
National Police Agency, Ministry of the Interior. 2024. {``165
Anti-Fraud Hotline.''}
\url{https://www.npa.gov.tw/en/app/artwebsite/view?id=8035\&module=artwebsite\&serno=ed2427e1-de0a-4f6f-8f68-8f83b604e89b}.

\leavevmode\vadjust pre{\hypertarget{ref-rugayan2023asd}{}}%
Rugayan, Janine, Giampiero Salvi, and Torbjørn Svendsen. 2023.
{``Perceptual and Task-Oriented Assessment of a Semantic Metric for ASR
Evaluation.''} In \emph{Interspeech 2023}, 2158--62.
\url{https://doi.org/10.21437/Interspeech.2023-1778}.

\end{CSLReferences}

\end{document}
